% Fichier généré automatiquement par tools/generate_corrections.py.
% Source : RDM/RDM-05-Deformation/541_RdM/541_RdM.tex
% Statut : complet (2/2 question(s) renseignée(s), 0 reprise(s)).
\begin{corrigebox}[Corrigé extrait de la banque]
\CorrigeQuestion{1}
L'équation de la déformée est donnée par : $EI_{GZ}y''(x)=M_{fz}(x)$.

\textbf{Calcul du torseur de cohésion}

La poutre est composée de 2 tronçons :
\begin{itemize}
\item 1\ier tronçon : $\lambda\in[0,L]$
\item 2\ieme tronçon : $\lambda\in[L,2L]$
\end{itemize}

\textbf{1\ier tronçon}
\begin{itemize}
\item On isole la partie II.
\item BAME : 
\begin{itemize}
\item $\torseurstat{T}{I}{II}$
\item $\torseurstat{T}{F}{II}$
\end{itemize}
\item D'après le PFS, on a donc $\torseurstat{T}{I}{II} + \torseurstat{T}{F}{II} = \{0\}$ et donc 
$\torseurstat{T}{II}{I} =  \torseurstat{T}{F}{II}$.
\end{itemize}
On a donc $\forall \lambda\in[0,L]$, $\vectm{G}{F}{II}=  \left(L-\lambda\right)\vect{x} \wedge - F\vect{y} =-F\left(L-\lambda\right) \vect{z} $. 

Au final, $\torseurstat{T}{II}{I} = \torseurl{-F\vect{y}}{-F\left(L-\lambda\right) \vect{z}}{G}$.

\textbf{2\ieme tronçon}

$\forall \lambda \in [0,L]$, $\torseurstat{T}{II}{I} = \torseurl{\vect{0}}{\vect{0}}{G}$.

\textbf{Calcul de la déformée sur le premier tronçon}

 $\forall x \in [L,2L]$, $EI_{GZ}y''(x)=-F\left(L-x\right) =-FL + Fx $ 
et en intégrant $EI_{GZ}y'(x)= -FLx +\dfrac{1}{2}Fx^2+ c_1$ et  $EI_{GZ}y(x)= -\dfrac{1}{2}FLx^2 +F \dfrac{1}{6}x^3+ c_1x + c_2$.

La liaison en $x=0$ étant une encastrement, on a $y(0)=0$ et $y'(0)=0$. 
En conséquence, $c_2=0 $ et $c_1=0$.

Au final, $EI_{GZ}y(x)= -\dfrac{1}{2}FLx^2 + \dfrac{1}{6}Fx^3 = Fx^2\dfrac{x - 3L}{6}$.

On peut alors exprimer la flèche en $L$ : $EI_{GZ}y(L)= -\dfrac{ FL^3}{3}$.

\textbf{Calcul de la déformée sur le second tronçon}

$\forall x \in [L,2L]$, $EI_{GZ}y_2''(x)=0$ et en intégrant $EI_{GZ}y_2'(x)=c_3$ et $EI_{GZ}y_2(x)=c_3x + c_4$. 
On a de plus $y(L) = y_2(L)$ et $y'(L) = y_2'(L)$. 

D'une part, $F\dfrac{3L^2- 6L^2}{6}=c_3 $ et donc $c_3 =  -F\dfrac{ L^2}{2} $

D'autre part, $FL^2\dfrac{L - 3L}{6}=-F\dfrac{ L^2}{2} L + c_4$ et $c_4 =  FL^2\dfrac{L - 3L}{6}+F\dfrac{ L^2}{2} L  = \dfrac{FL^3}{6}$.
\CorrigeQuestion{2}
\par\medskip\begin{center}\captionsetup{type=figure}
\centering

\end{center}\par\medskip
\end{corrigebox}
